\chapter{対称性、縮退、角運動量のはしご}
\providecommand{\ket}[1]{\left\lvert#1\right\rangle}

\section{対称性と縮退したエネルギー準位}

対称性とは、系の物理的記述を変えない操作である。それはハミルトニアン、Schr\"odinger方程式、あるいは取り得るエネルギーを変えない。ユニタリ演算子\(S\)がその操作を表すならば、不変性とは
\begin{equation}
S H S^\dagger=H,
\end{equation}
を意味し、あるいは等価的に
\begin{equation}
[S,H]=0.
\end{equation}
である。

並進はその一例である。原子を一つの場所から別の場所へ移しても、そのエネルギーは変わらない。これは基底状態だけでなく、すべての原子状態について成り立つ。回転もまた一例であり、ここでの主要な主題となる。

二つの同一粒子を入れ替えることもまた対称性である。二つの電子を電子1と電子2、あるいはハリーとフレッドと呼んでも、それらが異なるものになるわけではない。ラベルを入れ替えても原子の性質は変わり得ない。

結晶にはさらに別種の並進対称性がある。結晶全体を空間中で動かすのは通常の並進である。しかし無限結晶に対しては、一格子間隔だけ並進させても同じ物理的配置が得られる。等価的に、すべての格子点の指数を1だけずらすことによって格子点を付け替えることができる。境界は有限結晶に対してこの厳密な主張を損なうが、非常に長い結晶の端から遠く離れたところでは良い対称性となる。

空間の一様性---すべての位置が他のすべての位置と同じであるという事実---は並進対称性である。空間の等方性---すべての方向が他のすべての方向と同じであるという事実---は回転対称性である。

縮退は対称性と関係しているが、同じものではない。あるエネルギー準位は、複数の状態がちょうどそのエネルギーをもつとき縮退しているという。離散スペクトルをもつ一般的なハミルトニアンに対して、無関係な二つの準位の間の厳密な等しさは驚くべき偶然であろう。\(10^{15}\)分の1までの一致であっても、依然として厳密な等しさではない。厳密な縮退は通常、対称性を指し示す。

対称性はときに縮退を意味するが、常にそうではない。最初の例は、なぜそうなり得るのか、そしてなぜ一つの対称性だけでは不十分であり得るのかの両方を示す。

\section{円周上の粒子}

円形の針金の上を運動するよう拘束されたビーズを考える。その位置はデカルト座標\(x\)と\(y\)で記述することもできるが、一つの角度座標で十分である：
\begin{equation}
\psi=\psi(\theta).
\end{equation}
円周上の位置確率密度は\(\psi^*(\theta)\psi(\theta)\)である。

小さな角度\(\epsilon\)だけ反時計回りに能動的に回転させると、
\begin{equation}
\psi(\theta)\longrightarrow\psi(\theta-\epsilon).
\end{equation}
したがって
\begin{align}
\delta\psi(\theta)
&=\psi(\theta-\epsilon)-\psi(\theta)\\
&=-\epsilon\frac{d\psi}{d\theta}
+\mathcal O(\epsilon^2).
\end{align}

演算子
\begin{equation}
L_z=-i\hbar\frac{d}{d\theta}
\end{equation}
は周期的な波動関数上でエルミートである。これは
\begin{equation}
p_x=-i\hbar\frac{d}{dx}.
\end{equation}
の角度版である。\(i\)の因子はエルミート性にとって不可欠である。\(L_z\)を用いると、
\begin{equation}
\delta\psi
=-\frac{i\epsilon}{\hbar}L_z\psi
+\mathcal O(\epsilon^2),
\end{equation}
であるから、無限小回転演算子は
\begin{equation}
R(\epsilon)
=1-\frac{i\epsilon}{\hbar}L_z
+\mathcal O(\epsilon^2).
\end{equation}
となる。これにより\(L_z\)が円周に垂直な軸のまわりの回転の生成子であることが同定される。

\begin{classroomqa}
\lecturetimestamp{00:11:19}
\audiencequestion{\(i\)を含む因子の一方は逆符号であるべきではないか？}
\lecturerresponse{符号を確認すると上の関係が得られる。規約\(\psi(\theta)\to\psi(\theta-\epsilon)\)では、変化は\(-\epsilon\,d\psi/d\theta\)であり、これは\(-i\epsilon L_z\psi/\hbar\)に等しい。}
\end{classroomqa}

次に角運動量の固有状態を求める：
\begin{equation}
L_z\psi=M\psi.
\end{equation}
微分演算子を代入すると
\begin{equation}
-i\hbar\frac{d\psi}{d\theta}=M\psi,
\end{equation}
が得られ、その解は
\begin{equation}
\psi(\theta)=C\exp\left(\frac{iM\theta}{\hbar}\right).
\end{equation}
の形をとる。

無限直線上の運動には現れない追加の条件がある。点\(\theta=0\)と\(\theta=2\pi\)は同じ点であるから、一価の軌道波動関数は
\begin{equation}
\psi(\theta+2\pi)=\psi(\theta).
\end{equation}
を満たさなければならない。上の固有関数に対して、これは
\begin{equation}
\exp\left(\frac{2\pi iM}{\hbar}\right)=1.
\end{equation}
を要求する。したがって、
\begin{equation}
M=m\hbar,
\qquad
m\in\mathbb Z.
\end{equation}
円周上で可能な軌道角運動量はしたがって
\begin{equation}
\ldots,-2\hbar,-\hbar,0,\hbar,2\hbar,\ldots.
\end{equation}
である。

角運動量\(m\hbar\)をもつ状態は、角運動量\(-m\hbar\)をもつ状態と同じエネルギーをもたなければならないのだろうか？正と負の\(m\)は反対向きの循環を記述する。問題は、回転対称性だけがそれら二つの運動に等しいエネルギーを要求するかどうかである。

\begin{classroomqa}
\lecturetimestamp{00:22:51}
\audiencequestion{\(\psi(\theta-\epsilon)\)を用いることは、粒子が時計回りに運動していることを意味するのか？}
\lecturerresponse{いいえ。その規約はどの方向を正の回転と呼ぶかを定めるだけである。それは角度位置と角運動量の符号を定義するが、粒子が実際に運動する方向を決定するものではない。}
\end{classroomqa}

円周の面に垂直な磁場が反例を与える。それは正と負の角運動量のエネルギーを反対方向にずらし、それによってその対を分裂させ得る。ずれの符号は電荷と配向の規約に依存するが、分裂が重要な点である。

そのような場は依然として垂直軸のまわりの回転を保存する。したがって、回転対称性それ自体は
\begin{equation}
E(m)=E(-m).
\end{equation}
を要求しないことを示している。分裂がなければ、\(m=1\) と \(m=-1\) は縮退した対をなし得るし、\(m=2\) と \(m=-2\) も同様であり得る。状態 \(m=0\) には、この種の異なる相手が存在しない。

\section{反射対称性と \texorpdfstring{\(m,-m\)}{m, -m} 対}

離散的な反射対称性を加える。\(\mathcal M\) を円の直径に関する反射を表すものとする：
\begin{equation}
(\mathcal M\psi)(\theta)=\psi(-\theta).
\end{equation}
それは循環の向きを逆にする：
\begin{equation}
\mathcal M\psi_m=\psi_{-m}.
\end{equation}

反射が対称性ならば、
\begin{equation}
[\mathcal M,H]=0.
\end{equation}
である。
\begin{equation}
H\ket{m}=E_m\ket{m}.
\end{equation}
と仮定する。\(\mathcal M\) を適用すると
\begin{align}
H\mathcal M\ket{m}
&=\mathcal M H\ket{m}\\
&=E_m\mathcal M\ket{m}.
\end{align}
が得られる。\(\mathcal M\ket{m}\) は角運動量 \(-m\hbar\) をもつ状態であるから、
\begin{equation}
E_{-m}=E_m.
\end{equation}
である。回転対称性は角運動量を保存し、反射対称性は二つの反対の値を関係づける。両者を合わせると、その対が縮退していることが要求される。

\begin{classroomqa}
\lecturetimestamp{00:29:54}
\audiencequestion{\(m\) モードと \(-m\) モードは同じ確率密度をもつので、それだけでそれらのエネルギーが等しいことが示されるのか？}
\lecturerresponse{いいえ。波動関数 \(e^{im\theta}\) と \(e^{-im\theta}\) は、その絶対値の二乗が一致していても異なる。角運動量固有値方程式は波動関数を関与させ、\(\psi^*\psi\) だけではない。量子力学は確率だけでなく振幅を要求する。}
\end{classroomqa}

磁場はその軸のまわりの回転を保存するが、上で用いた反射を破る。

\begin{classroomqa}
\lecturetimestamp{00:33:05}
\audiencequestion{反射対称性は磁場を逆にするのか？}
\lecturerresponse{はい。磁場は軸性ベクトルである。板に垂直に入る向きの場を反射すると、板から出る向きの場が得られる。等価に、電流ループとその鏡像は逆向きの電流を運び、逆向きの場を生じる。}
\end{classroomqa}

回転と反射の非可換性は直接に見ることができる。角運動量固有関数に対して、
\begin{align}
\mathcal M L_z\psi_m
&=m\hbar\,\psi_{-m},\\
L_z\mathcal M\psi_m
&=-m\hbar\,\psi_{-m}.
\end{align}
したがって
\begin{equation}
[\mathcal M,L_z]\psi_m
=2m\hbar\,\psi_{-m},
\end{equation}
であり、これは \(m\neq0\) に対して非零である。回転してから反射する操作は、反射してから回転する操作と同じではない。

両方の変換が対称性であるとき、両方の演算子は \(H\) と可換であるが、互いに可換ではない。これが縮退の原因となるパターンである。

\begin{classroomqa}
\lecturetimestamp{00:40:28}
\audiencequestion{\(x\) と \(p\) のような非可換なハイゼンベルク対は縮退を生じるのか？}
\lecturerresponse{一般にはそうではない。演算子はまた対称性でなければならない：それらはハミルトニアンと可換でなければならない。位置と運動量は通常、両方が保存されるわけではない。もし両方が \(H\) と可換ならば、それらの非可換性は縮退を意味するであろう。}
\end{classroomqa}

\section{対称性の交換子代数}

\(A\) と \(B\) が二つの連続対称性のエルミート生成子であるとする：
\begin{equation}
[A,H]=0,
\qquad
[B,H]=0.
\end{equation}
これらの方程式のどこにも、\(A\) と \(B\) が互いに可換であることは要求されていない。

交換子 \([A,B]\) は反エルミートであるから、エルミート演算子
\begin{equation}
C=i[A,B].
\end{equation}
を定義する。数因子と全体の符号はここでは重要でない。重要な点は、\(C\) もまたハミルトニアンと可換であることである。実際、
\begin{equation}
[AB,H]
=A[B,H]+[A,H]B=0,
\end{equation}
であり、同様に
\begin{equation}
[BA,H]=0.
\end{equation}
である。したがって
\begin{equation}
[C,H]
=i\bigl([AB,H]-[BA,H]\bigr)=0.
\end{equation}

新しい生成子 \(C\) は \(A\) と \(B\) から独立であり得るし、すでに知られている生成子の線形結合にすぎないこともあり得る。もし新しいならば、それをリストに加えてさらなる交換子を作る。この過程を繰り返すと無限に続くこともあるが、しばしば有限個の生成子で閉じる：それ以上の交換子はすべて、すでに存在するものの線形結合である。その閉じた系が交換子代数またはリー代数である。

有限の対称変換は群をなす。それらの無限小生成子は対応するリー代数をなす。

縮退が重要であるのは、それがスペクトルを組織するからである。二つの状態が等しいエネルギーをもつならば、一方がそのエネルギー差を運ぶ光子を放出して他方へ崩壊することはできない。原子分光学では、準位の配置と縮退が可能なスペクトル線を決定する。同じ論理が励起ハドロンとその遷移にも当てはまる。

\begin{classroomqa}
\lecturetimestamp{00:49:07}
\audiencequestion{これらの交換子自体は対称群の要素なのか？}
\lecturerresponse{ここで議論している演算子は生成子である。生成子とそれらの交換子は交換子代数、すなわちリー代数をなす。生成子から作られる有限の対称操作が群の要素である。}
\end{classroomqa}

すべての対称性生成子が互いに可換ならば、その対称性は可換であると呼ばれる。それらの交換子のいくつかが非零ならば、非可換である。

\begin{classroomqa}
\lecturetimestamp{00:50:25}
\audiencequestion{生成子の基本的な定義とは何か。}
\lecturerresponse{恒等変換に非常に近いユニタリ対称変換をとり、それを\(1-i\epsilon G/\hbar+\cdots\)と書く。エルミート演算子\(G\)がその生成子である。対応する無限小変換を繰り返せば有限変換が構成される。}
\end{classroomqa}

\begin{classroomqa}
\lecturetimestamp{00:52:08}
\audiencequestion{対称性の生成子の代数とは何を意味するのか。}
\lecturerresponse{それは、加算、減算、通常の数による乗算、および交換ができる生成子の集合である。この代数では交換子が乗算の役割を果たし、閉性とはこれらの演算が集合の外の生成子を生じないことを意味する。}
\end{classroomqa}

\section{異なる軸のまわりの回転}

回転は非可換対称性の最も重要な例を与える。\(x\)軸のまわりの回転に続いて\(y\)軸のまわりの回転を行っても、一般には逆の順序で行った場合と同じにはならない。それらの不一致自体が第三の軸のまわりの回転と関係している。

\begin{classroomqa}
\lecturetimestamp{00:54:22}
\audiencequestion{有限回転は可換でないのに、無限小回転は可換なのか。}
\lecturerresponse{それらは二つの小さな回転角について二次の差がある。それで十分である。すなわち、それらの小さなパラメータをくくり出した後、生成子自身が非零の交換子をもつ。}
\end{classroomqa}

量子交換子を計算する前に、対応する古典的描像を考える。角運動量がその平面に垂直な円形の惑星軌道をとる。軌道全体を\(y\)軸のまわりに回転させると、全く同じエネルギーをもつ別の許された軌道が得られる。しかしその角運動量ベクトルは今や異なる\(x\)成分をもつ。

したがって回転不変性は、異なる角運動量成分をもつ同じエネルギーの配置を生み出す。量子力学的には、これは異なる軸のまわりの回転の生成子の間の非可換性として現れる。

\section{デカルト座標における角運動量}

まず、粒子を円周上に制限することなく、\(xy\)平面内を運動するものとして考える。微小な反時計回りの回転の下で、
\begin{equation}
\delta x=-\epsilon y,
\qquad
\delta y=\epsilon x.
\end{equation}
能動的に回転させた状態に対して、変換された波動関数は逆回転した座標で評価される：
\begin{equation}
\psi_\epsilon(x,y)
=\psi(x+\epsilon y,y-\epsilon x).
\end{equation}
したがって、
\begin{align}
\delta\psi
&=\epsilon\left(
y\frac{\partial}{\partial x}
-x\frac{\partial}{\partial y}
\right)\psi
+\mathcal O(\epsilon^2)\\
&=-\frac{i\epsilon}{\hbar}
\left(xp_y-yp_x\right)\psi
+\mathcal O(\epsilon^2).
\end{align}
よって、\(z\)軸のまわりの回転の生成子は
\begin{equation}
L_z=xp_y-yp_x.
\end{equation}
である。これは、よく知られた古典的ベクトル
\begin{equation}
\mathbf L=\mathbf r\times\mathbf p.
\end{equation}
の\(z\)成分である。それは力学的角運動量だけではなく、量子力学においては波動関数の回転を生成する演算子である。\footnote{能動的回転の規約は、先の変換\(\psi(\theta)\mapsto\psi(\theta-\epsilon)\)および、口頭での導出における途中の符号の不確定性の後の最終的な同定\(L_z=xp_y-yp_x\)と一致するように選ばれている。}

\(x\to y\to z\to x\)を巡回させると他の成分が得られる。\footnote{\(L_x\)および\(L_y\)の式は、講義で明示的に規定された巡回置換\(x\to y\to z\to x\)を\(L_z=xp_y-yp_x\)に適用することによって再構成される。}
\begin{align}
L_x&=yp_z-zp_y,\\
L_y&=zp_x-xp_z,\\
L_z&=xp_y-yp_x.
\end{align}

回転不変な系に対して、これらの演算子はハミルトニアンと交換する。
\begin{equation}
[H,L_x]=[H,L_y]=[H,L_z]=0.
\end{equation}
中心力場中の粒子が標準的な例である。これらの式は量子形式における角運動量保存である。

\section{角運動量リー代数}

正準交換関係は
\begin{align}
[x_i,x_j]&=0,\\
[p_i,p_j]&=0,\\
[x_i,p_j]&=i\hbar\,\delta_{ij}.
\end{align}
座標は異なる方向に属する運動量成分と交換する。したがって例えば\(x\)は\(p_y\)と交換するが、\(x\)は\(p_x\)と交換しない。

次に
\begin{equation}
[L_x,L_y]
=
[yp_z-zp_y,\;zp_x-xp_z].
\end{equation}
を計算する。ほとんどの項は消える。非零の二つの寄与は
\begin{align}
[yp_z,zp_x]&=-i\hbar\,yp_x,\\
[zp_y,xp_z]&=i\hbar\,xp_y.
\end{align}
である。したがって
\begin{align}
[L_x,L_y]
&=i\hbar\left(xp_y-yp_x\right)\\
&=i\hbar L_z.
\end{align}
座標ラベルを巡回させると
\begin{align}
[L_x,L_y]&=i\hbar L_z,\\
[L_y,L_z]&=i\hbar L_x,\\
[L_z,L_x]&=i\hbar L_y.
\end{align}

\begin{classroomqa}
\lecturetimestamp{01:14:50}
\audiencequestion{交換子が\(L_z\)を生み出すとき、\(L_z\)が\(L_x\)と\(L_y\)の線形結合でないことを依然として示さなければならないのか。}
\lecturerresponse{はい。その独立性を仮定するのではなく確立する義務がある。この時点では直接的な代数的証明は行わない。}
\end{classroomqa}

表示された交換子は\(L_x\)、\(L_y\)、\(L_z\)の張る空間上で閉じている。すなわち、これら三つをすべて含めれば、それ以上の交換子はその空間の外の演算子を生じない。これが角運動量リー代数である。
\begin{equation}
[H,L_i]=0,
\qquad i=x,y,z,
\end{equation}
と合わせて、これらの関係は、角Schr\"odinger方程式を直接解くことを要せずに、エネルギー準位についての相当な情報を含んでいる。

同様の代数的戦略は量子調和振動子にも用いられ、そこでは生成演算子と消滅演算子がそのスペクトルを生成する。ここでは、類似の演算子はエネルギーではなく角運動量の \(z\) 成分を変化させる。

\(z\) 軸を選ぶことは表現の選択にすぎない。\(z\) が \(x\) や \(y\) より物理的に好ましいわけではないが、一つの成分を選ぶと便利な基底が得られる。
\begin{equation}
L_z\ket{m}=m\hbar\ket{m}.
\end{equation}
課題は、\(m\) のどの値が起こり得るか、そして異なる \(m\) をもつ状態がどのように関係するかを決定することである。

\begin{classroomqa}
\lecturetimestamp{01:24:02}
\audiencequestion{この構成を調和振動子と比較するというのは、量子調和振動子のことですか。それは前回の授業では扱われませんでした。}
\lecturerresponse{はい、比較は量子振動子との比較である。現在の角運動量の議論は自己完結しており、振動子の知識を必要としない。この類推は後に逆方向にも読める：振動子は同様の代数学的方法を用いる。}
\end{classroomqa}

\section{角運動量の昇降}

定義する
\begin{equation}
L_+=L_x+iL_y,
\qquad
L_-=L_x-iL_y.
\end{equation}
角運動量の交換関係は次を意味する
\begin{equation}
[L_z,L_+]=\hbar L_+,
\qquad
[L_z,L_-]=-\hbar L_-.
\end{equation}

最初の関係を \(L_z\) の固有状態に適用する：
\begin{align}
L_zL_+\ket{m}
&=\left(L_+L_z+[L_z,L_+]\right)\ket{m}\\
&=(m+1)\hbar\,L_+\ket{m}.
\end{align}
\(L_+\ket{m}=0\) でない限り、それは \(L_z\) の別の固有状態であり、今度は固有値 \((m+1)\hbar\) をもつ。同様に、
\begin{equation}
L_zL_-\ket{m}
=(m-1)\hbar\,L_-\ket{m}.
\end{equation}
したがって \(L_+\) は \(m\) を 1 だけ上げ、\(L_-\) はそれを 1 だけ下げる。

任意の一つの固有値から出発して、代数は整数で隔てられた列を生成する。その列は無限に続くこともあれば、上端で
\begin{equation}
L_+\ket{m_{\mathrm{top}}}=0
\end{equation}
下端で
\begin{equation}
L_-\ket{m_{\mathrm{bottom}}}=0
\end{equation}
となるために終わることもある。

\begin{classroomqa}
\lecturetimestamp{01:35:11}
\audiencequestion{角運動量は円の議論によってすでに整数であることが要求されていたのではないですか。}
\lecturerresponse{その議論は、一価の位置波動関数で表される軌道角運動量に対して整数値を証明する。抽象的な角運動量代数はより緩い結果を与える：はしごは整数値または半整数値のいずれかを含み得る。}
\end{classroomqa}

はしごに端がある場合、回転不変性によってその両端が関係づけられる。\(z\)に垂直な軸のまわりに\(180^\circ\)回転すると、次の変換が生じる。
\begin{equation}
L_z\longrightarrow-L_z.
\end{equation}
したがって、スペクトルはゼロを中心に対称でなければならない。
\begin{equation}
m_{\mathrm{bottom}}=-m_{\mathrm{top}}.
\end{equation}
隣り合う値の差は1なので、可能性は二つである。はしごに含まれる値は整数
\begin{equation}
\ldots,-2,-1,0,1,2,\ldots,
\end{equation}
または半整数
\begin{equation}
\ldots,-\frac{5}{2},-\frac{3}{2},-\frac{1}{2},
\frac{1}{2},\frac{3}{2},\frac{5}{2},\ldots,
\end{equation}
であり、昇演算子または降演算子を作用させてゼロになるところが端点となる。

円上の一価波動関数の議論は、軌道角運動量に対して整数の可能性を選ぶ：
\begin{equation}
\mathbf L=\mathbf r\times\mathbf p.
\end{equation}
半整数の可能性も物理的であり、スピン角運動量に属し、それは別に扱われる。

最後に、ハミルトニアンが回転不変であるとき、はしごの状態は等しいエネルギーをもつ。仮定する
\begin{equation}
H\ket{m}=E\ket{m}.
\end{equation}
各成分 \(L_i\) は \(H\) と交換するので、\(L_+\) と \(L_-\) も同様である。したがって
\begin{align}
H L_+\ket{m}
&=L_+H\ket{m}\\
&=E L_+\ket{m}.
\end{align}
\(L_+\ket{m}\neq0\) のときはいつでも、上げられた状態は元の状態と同じエネルギーをもつ。同じ計算から
\begin{equation}
H L_-\ket{m}=E L_-\ket{m}.
\end{equation}
が得られる。したがって昇降は、異なる \(L_z\) 固有値をもつがまったく同じエネルギーをもつ状態の族を生成する。これが、講義全体で展開された非可換対称性と縮退の間の関係である。

\begin{classroomqa}
\lecturetimestamp{01:44:13}
\audiencequestion{これらのはしごの状態は、異なる回転軸を選ぶことと関係していますか。}
\lecturerresponse{それらは回転によって関係づけられる。\(L_x\) または \(L_y\) の固有ベクトルは、選ばれた \(L_z\) 固有ベクトルの線形結合である。状態を回転させると、そのエネルギーを変えずに角運動量成分を変える。}
\end{classroomqa}

\(L_+\)と\(L_-\)の選択は、これらの組み合わせがはしごを直接露わにするので便利である。この技巧は量子力学の至るところで再現する。交換子代数が閉じるとき、昇降演算子はしばしばその固有値スペクトルを組織化する。

\begin{classroomqa}
\lecturetimestamp{01:46:33}
\audiencequestion{\(L_+\)が最大状態に作用すると何が起こるか、そして頂上はどのように定義されるのか。}
\lecturerresponse{それはゼロを与える。それが終端するはしごの頂上の定義である。同様に\(L_-\)は底でゼロを与える。抽象的なはしごは終端する必要はない。後に考察する原子の角運動量状態に対しては、通常、終端するはしごに出会う。}
\end{classroomqa}

次の段階は、この角運動量の構造をさらに発展させ、水素を含む原子に適用することである。ここで見出された回転多重項は縮退パターンの一部を説明する。さらなる構造はこれから展開されるべきものである。
